\section{Discussion}

The ledger, in one place. Completion guarantees cost a median 10.3
percent of goodput per provisioned GPU-hour in the tested region, and
the cost is U-shaped in load: deepest, 22 to 24 percent, exactly at
nominal provisioning, recovering toward saturation, inverting past it.
They buy three measured things: interactive tail-latency isolation of 63 to 97 percent across most of the region, with one identified inversion
corner; containment of adversarial runaways to a 3.1-point degradation
where the strongest request-oriented schedulers starve compliant agentic
work to zero; and a regime boundary at 1.5x load beyond which the
guarantee begins to pay for itself. The efficiency thesis, that the
purchases outweigh the cost on a single hardware-normalized metric, is
dead on its own pre-registered terms. What survives is a price list.

The verdict closes each of Section 3's transfers on its own terms.
Lesson one's economics largely do not survive the crossing: the field
that invented per-flow guarantees abandoned them substantially for their
cost, and this evaluation measured the same abandonment logic operating
at the serving layer, with the added force that the tuner, offered a
range of reserved-subset sizes under an honest metric, chose the
smallest in every workload family at every horizon. The mechanism,
optimized truthfully, minimizes itself; the duty-cycle analysis of
Section 7.5 says why, and says it is structural rather than
configurational. Lesson two's shape survives but its transferred signal
does not: degradation should still act early, probabilistically, and per
class, but keying it to committed headroom rather than queue state,
this design's one deliberate departure from the RED lineage, delivered
no measurable advantage at the system level (Section 7.3) and no
attributable one at the component level (Section 7.5), and the forward
signal this paper argued the serving layer gets for free is, on this
evidence, worth nothing that queue depth does not already provide.
Lessons three and four cross intact, because they were never economic
claims: fate-independent supervision and provenance-bound admission are
authority structures, their cost is engineering rather than goodput, and
the one property of theirs this evaluation could measure, the budget's
containment of adversarial runaways, held exactly as specified while
the strongest priority schedulers starved every compliant flow they
were nominally serving. The contract-versus-priority distinction,
argued in Section 4.2 as a matter of kind, appears in Section 7.4 as a
matter of record.

The honest reading of the mechanism, then, is not ``efficiency technique
that failed'' but ``insurance product whose premium is now known.'' Fleets
that are well-provisioned and run nominal load sit at the bottom of the
U and should not reserve: the premium is maximal there, and the
starvation pathology, while real, is avoidable with cheaper instruments
at low agentic share. The case for reservations concentrates where the
purchases concentrate: fleets that price interactive tail latency highly
and currently buy isolation through overprovisioning; fleets with
adversarial or untrusted agentic exposure, where containment is not
replicable by per-request caps; and fleets that run hot, past
saturation, where the sign has already flipped. Even there, the
K-to-floor result disciplines enthusiasm: reserve for the few flows
whose completion is worth a known premium, and for nothing else, which
is what Section 3's first lesson prescribed before the evaluation
priced it.

Three things could move the verdict, and we name them rather than imply
them. Hardware calibration: the binding sensitivity of Section 7.6
means the deficit's magnitude, not its sign, is uncertain at roughly
the uncertainty of one anchored constant; a calibrated rerun is
mechanical if the constant is ever measured. Workload shape: gap
structure drives stranding, and agentic workloads with tighter
tool-time distributions than our anchors would shrink the premium.
Objective: this evaluation priced everything into one number by design;
an operator who prices tail latency or adversarial robustness as
first-class objectives is solving a different optimization, and Section
7's decomposition is written to let them price it. What cannot move the
verdict is re-measurement under this registration: the one-fix clause
spent its fix, and the number stands.
